\newpage

\section{Предлагаемый метод анализа распределения параметров нейросети}
В качестве априорного распределения параметров $\mathbf{w}$ рассмотрим $k$-мерное распределение Лапласа $\Lap$ с плотностью распределения:
\begin{equation} \label{methods.1}
f(\mathbf{x}) = \frac{2e^{\mathbf{x}^{\mathsf{T}}\mbox{\boldmath$\Sigma^{-1}$}\mbox{\boldmath$\mu$}}}{(2\pi)^{\frac{k}{2}}|\mbox{\boldmath$\Sigma$}|^{0,5}} \left( \frac{\mathbf{x}^{\mathsf{T}}\mbox{\boldmath$\Sigma^{-1}$}\mathbf{x}}{2 + \mbox{\boldmath$\mu$}^{\mathsf{T}}\mbox{\boldmath$\Sigma^{-1}$}\mbox{\boldmath$\mu$}} \right)^{\frac{v}{2}} K_v \left( \sqrt{(2 + \mbox{\boldmath$\mu$}^{\mathsf{T}}\mbox{\boldmath$\Sigma^{-1}$}\mbox{\boldmath$\mu$})(\mathbf{x}^{\mathsf{T}}\mbox{\boldmath$\Sigma^{-1}$}\mathbf{x})} \right),
\end{equation}
где $\mbox{\boldmath$\mu$} \in \mathbb{R}^k$ -- математичекое ожидание, $\mbox{\boldmath$\Sigma$} \in \mathbb{R}^{k \times k}$ -- матрица ковариации, $v = (2-k)/2$ и $K_v$ -- модифицированная функция Бесселя второго рода.

Рассмотрим случай, когда $k = 2$, $\mathbf{m} = \mathbf{0}$, а матрицы $\mbox{\boldmath$\Theta$}$ и $\mbox{\boldmath$\Sigma$}$ диаганальны. В таком случае плотность распределения $p(\mathbf{w} \mid \mathcal{A}) \sim \Lap(\mathbf{m}, \mbox{\boldmath$\Theta$}) \sim \Lap(\mathbf{0}, \mbox{\boldmath$\Theta$})$ будет иметь вид:
\begin{equation} \label{methods.2}
f(\mathbf{x}) = \frac{2}{(2\pi)^{\frac{k}{2}}|\mbox{\boldmath$\Theta$}|^{0,5}} K_0 \left( \sqrt{2 \mathbf{x}^{\mathsf{T}}\mbox{\boldmath$\Theta^{-1}$}\mathbf{x}} \right).
\end{equation}

Рассмотрим несколько случаев:
\begin{enumerate}
\item Пусть вектор параметров мал по норме, такой, что:
\begin{equation} \label{methods.3}
||\mathbf{w}|| \lesssim \frac{2}{||\mbox{\boldmath$\Sigma$}||}.
\end{equation}
В таком случае можно использовать аппроксимацию функции Бесселя \cite{Temme1975}:
\begin{equation} \label{methods.4}
K_v(x) \simeq \frac{\Gamma(v)}{2} \left( \frac{2}{x} \right)^v \Rightarrow K_0(x) \sim \frac{\Gamma(0)}{2} \sim \const.
\end{equation}
Тогда можно оценить:
\[
\label{methods.5}
\begin{aligned}
D_{\KL}(q(\mathbf{w}) \parallel p(\mathbf{w} \mid \mathcal{A})) = \int_{\mathbf{w \in \mathbb{W}_{\mathcal{J}}}} q(\mathbf{w}) \log \frac{p(\mathbf{w} \mid \mathcal{A})}{q(\mathbf{w})} d \mathbf{w} = 
\end{aligned}
\]

\[
\begin{aligned}
= \int_{\mathbf{w \in \mathbb{W}_{\mathcal{J}}}} q(\mathbf{w}) \log p(\mathbf{w} \mid \mathcal{A}) q(\mathbf{w}) d \mathbf{w} - \int_{\mathbf{w \in \mathbb{W}_{\mathcal{J}}}} q(\mathbf{w}) \log q(\mathbf{w}) d \mathbf{w} \sim
\end{aligned}
\]

\[
\begin{aligned}
\sim \int_{\mathbf{w \in \mathbb{W}_{\mathcal{J}}}} \frac{e^{\mathbf{w}^{\mathsf{T}}\mbox{\boldmath$\Sigma^{-1}$}\mbox{\boldmath$\mu$}}}{|\mbox{\boldmath$\Sigma$}|^{0,5}} K_0 \left( \sqrt{(2 + \mbox{\boldmath$\mu$}^{\mathsf{T}}\mbox{\boldmath$\Sigma^{-1}$}\mbox{\boldmath$\mu$})(\mathbf{w}^{\mathsf{T}}\mbox{\boldmath$\Sigma^{-1}$}\mathbf{w})} \right) \log \frac{1}{|\mbox{\boldmath$\Sigma$}|^{0,5}} \left(  K_0 \left( \sqrt{\mathbf{w}^{\mathsf{T}}\mbox{\boldmath$\Theta^{-1}$}\mathbf{w}} \right) \right) d \mathbf{w} -
\end{aligned}
\]

\[
\begin{aligned}
1
\end{aligned}
\]

\item Пусть норма вектора параметров большая:
\begin{equation} \label{methods.3}
||\mathbf{w}|| \sim ?????.
\end{equation}
Тогда функцию Бесселя можно аппроксимировать следующим образом [ссылка на https://dlmf.nist.gov/10.40]:
\begin{equation} \label{methods.3}
K_v(x) \simeq \left( \frac{\pi}{2x}\right)^{\frac{1}{2}}e^{-x} \sum \limits_{n=0}^{\infty} \frac{a_n(v)}{x^n} \Rightarrow K_0(x) \sim \frac{e^{-x}}{\sqrt{x}} \sum \limits_{n=0}^{10} \frac{a_n(0)}{x^n}.
\end{equation}
В этом случае можно получить следующую оценку:
\[1\] 
\end{enumerate}


